\scanpage{113}
\mspage{104}
% archive sheet number: 113

\[
\operatorname{supp}\underset{\substack{\parallel\\\Phi_{T,\alpha}}}{\unsure{\Phi|\varphi}}
=\operatorname{Omb}(\widetilde S_{ab})\setminus\alpha
=\operatorname{Omb}^{\circ}(S_{ab})\cup\beta
\]
\hfill(réunion disjointe)

On a déjà fait \textcircled{a} plus haut, il faut prouver encore \textcircled{\scriptsize $1^{\circ}$} et \textcircled{b}.

Prouvons $1^{\circ})$. Cela résulte, par réc. sur le nb de composantes connexes, du

% Tall marginal rule alongside the lemma statement, ending before the proof.
\noindent\begin{tikzpicture}[overlay]
\draw[line width=0.4pt] (-0.3,-0.45) .. controls (-0.34,-3.5) and (-0.22,-5.8) .. (-0.15,-6.3);
\end{tikzpicture}%
\uline{Lemme} Soit $X\in\unsure{\mathscr{M}|\mathcal{M}}$, et soient $\Phi,\Psi$ deux parties de $\mathscr{M}$ (préfigures ?) telles que
\[
\Phi\underset{\mathrm{pos}}{<}X<\Psi
\]
% margin note written slanted
\margin{Oui, il faut supposer $\Phi,\Psi$ des préfigures.}

[NB si $\Phi$ et $\Psi$ sont des préfigures, l'existence d'un tel $X$ équivaut à $\Phi\overset{\unsure{!|'}}{<}\Psi$ : si $\operatorname{ex}\delta\Phi\unsure{=|\leq}a<\operatorname{or}\delta\Psi=b$, on prend $X=S_{ab}$. Si $a=b$, on prend $X=P_a$].

Alors
\[
\operatorname{supp}^{\circ}(\Phi\cup\Psi)=\operatorname{supp}^{\circ}(\Phi)\cup\operatorname{supp}^{\circ}(\Psi)
\]

\uline{Démonstration} On a évidemment $\unsure{\supset|>}$, il faut prouver $\subset$, i.e.
\[
Z\in\operatorname{supp}^{\circ}(\Phi\cup\Psi)\overset{?}{\Longrightarrow}Z\in\operatorname{supp}^{\circ}\Phi\text{ ou }Z\in\operatorname{supp}^{\circ}\Psi.
\]
\deleted{\unsure{Car|On}} Notons que $\Phi\underset{\mathrm{pos}}{<}X\underset{\mathrm{pos}}{<}\Psi$ implique que
\[
X\in\operatorname{cosupp}^{\circ}(\Phi\cup\Psi),
\]
donc $Z\mathrel{|\circ|}X$, ce qui signifie aussi
\[
Z\underset{\mathrm{pos}}{<}X\text{ ou }Z\underset{\mathrm{pos}}{>}X.
\]
